\section{El ecosistema de modelos de código: de Copilot a la colaboración BigCode}

Una vez completo el pipeline de datos, la clase vuelve a la historia y al ecosistema de los code LLM. El cambio central es que Codex/Copilot demostró que la next-token prediction basta para constituir un baseline fuerte; después, la comunidad combinó datos, modelos y fine-tuning públicos en una gran cantidad de code assistants.

\subsection{2021: el código puede tratarse como texto con next-token prediction}

El code modeling temprano solía usar de forma explícita AST, program analysis o feature engineering; Codex mostró que, si se alimenta un Transformer grande con grandes cantidades de code, también puede aprender completion. Los métodos estructurados siguen teniendo valor, pero ya no son un requisito previo para llegar a una code generation fuerte.

\lecturefigure{slide-55.jpg}{GitHub Copilot 2021: el punto de partida como producto del next-token code modeling}{Deck oficial de Loubna Ben Allal, página 55}

\teachervoice{00:38:30--00:40:00, la ponente resume el significado histórico de Codex como «code can be treated like text»: no es que el código carezca de estructura, sino que un sequence baseline fuerte reduce mucho la barrera que imponían las características diseñadas a mano.}

\subsection{2024: más de 1.7K open code models y el panorama de los leaderboards}

En el Hub aparecen muchos modelos code-pretrained o con una code mixture; al mismo tiempo, los base models y los instruction-tuned models mejoran con rapidez en tablas de clasificación del tipo de HumanEval. La cantidad indica que el ecosystem está activo, pero no equivale a 1,700 pretraining runs independientes de alta calidad: muchos son fine-tunes, quantizations o derivatives.

\lecturefigure{slide-56.jpg}{Panorama de 2024: más de 1.7K open models trained on code}{Deck oficial de Loubna Ben Allal, página 56}

\lecturefigure{slide-57.jpg}{Panorama de las tablas de clasificación de code base models e instruction-tuned models fuertes}{Deck oficial de Loubna Ben Allal, página 57}

\begin{warningbox}{80\% en HumanEval no es «80\% de capacidad de programación»}
HumanEval consta de unas 164 Python function-completion tasks, evaluadas con unit tests. La puntuación depende del prompt, el sampling, pass@k, la contamination y el instruction tuning; no puede extrapolarse a repository editing, debugging, seguridad ni ingeniería de software con múltiples archivos.
\end{warningbox}

\subsection{Open code LLM landscape: provider, size y objective}

Meta Code Llama, BigCode StarCoder, DeepSeek Coder y otros forman distintas families; además están CodeGen, StableCode, CodeQwen y Granite. Al compararlos, conviene distinguir base/instruct, context, FIM, license, languages, training date y benchmark protocol.

\lecturefigure{slide-58.jpg}{Open code LLM landscape: Meta, BigCode, DeepSeek y otras families}{Deck oficial de Loubna Ben Allal, página 58}

\begin{knowledgebox}{Términos clave de los code models (primera aparición)}
Un \textbf{base model} aprende principalmente next-token/FIM; un \textbf{instruction-tuned model} aprende a responder instrucciones; \textbf{MQA} hace que todas las query heads compartan un solo conjunto de KV; \textbf{GQA} hace que varias query heads compartan un conjunto de KV, lo que por lo general es un punto intermedio entre calidad y KV-cache.
\end{knowledgebox}

\teachervoice{00:55:00--00:57:20, en la sesión de Q\&A la ponente dice que el entrenamiento con code y con text es similar en general, pero que un code assistant da más peso a long context, a la inference efficiency de MQA/GQA y a smaller models adecuados para el IDE (IDE-friendly).}

\subsection{BigCode: el objetivo de una colaboración de ciencia abierta}

BigCode no solo publica weights: hace que más de 1,000 investigadores e ingenieros participen en Slack y hace públicos el data processing, el training code, los models y una friendly license. El valor está en poder revisar el decision trail, no en tomar el collaboration size como garantía de calidad.

\lecturefigure{slide-59.jpg}{BigCode open-scientific collaboration: data, processing, models y comunidad transparentes}{Deck oficial de Loubna Ben Allal, página 59}

\lecturefigure{slide-60.jpg}{BigCode ecosystem: The Stack, StarCoder y los fine-tunes de la comunidad}{Deck oficial de Loubna Ben Allal, página 60}

\begin{importantbox}{La cadena de reutilización del ecosystem}
The Stack lo usan varios code models; a su vez, los checkpoints de StarCoder se vuelven la base de fine-tunes como StarChat y WizardCoder. La reutilización amplía el impacto, pero también significa que las decisiones de upstream sobre data quality, license y contamination se propagan hacia downstream.
\end{importantbox}

\teachervoice{00:40:00--00:43:10, la ponente dice que la razón directa para iniciar BigCode fue la full data transparency: la comunidad puede inspeccionar los datos (inspect data), reutilizar el processing code y revisar los modelos, en lugar de solo descargar weights.}

\subsection{Resumen de la sección}

El ecosistema de los code LLM pasó con rapidez del Copilot de código cerrado a datasets, base models y fine-tunes públicos. La siguiente pregunta ya no es «si hay modelos», sino si un release es verdaderamente abierto, transparente, reproducible y responsable.
